\section{合卷：西周留下的，不只是“分封制”三个字}

西周的问题从来不是抽象地选择“中央集权”还是“地方分权”。\eventref{WZ-1046-MUYE}{西周·牧野}之后，周王室面对的是一个比关中大得多、政治关系也更稠密的东方。王室既无条件在每处长期驻军，也不能让旧商中心自行恢复。于是，它把宗亲、功臣、旧有人群、城邑、礼器与军役拼接成网络；这不是预先画在纸上的制度蓝图，而是在每一次受命、安置和争端中才变得可用的关系。

\subsection{它解决了什么}

\eventref{WZ-1040-DONGZHENG}{西周·三监东征}后的成周、卫地与各处封国，使王命不必每次都从宗周出发。成周提供了会集、转运和册命的东方支点；地方家族则以自己的宗族、土地、车马和祭祀承接守卫、经营和联络。\eventref{WZ-EAST-ENFEOFFMENT}{西周·重组东方封国网络}不是把地方交出去便不再过问，而是让中心以较少的直接人手进入更广空间。王室仍须亲临、赏赐、裁断和调兵，封国也仍须在这些场合承认王命；正是这种往复，使网络在成康以后能够延展。

礼器让这套安排具有可见性。何尊、大盂鼎、册命簋和许多家族所作的铜器，都把一次命令、一次服务或一笔赏赐交给祖先和子孙。这样做解决的不是仪式的装饰问题：远方的受命者需要能证明自己为何有资格支配资源，王室也需要使功劳、责任与报偿有可反复调用的凭据。资源因此不只流向王室，也通过田宅、贝、车马、器物和名分流向中介家族。

\subsection{谁承担了代价}

一张网络若要运行，便不能只靠诰命。周原、丰镐与成周的仓储、道路、居址和作坊，提示粮食、铜料、工匠、役夫和运输必须被组织起来；裘卫器、卫器和散氏盘又显示贵族之间的资源关系也要反复交割、见证和申明。现存材料最常保存得到赏赐、能够作器的人，极少让耕作者、筑城者、车夫和久戍者以自己的名字出现。可以确定的是他们的劳动构成供给条件，不能确定的是每一群体实际负担多少赋役、又怎样评价这套秩序。

边疆把这种成本推到最尖锐处。昭王南行不返，说明距离、水路和地方关系可以使远征的风险骤然兑现；宣王时期反复出现的北方军旅、料民、戍守与人马管理，则显示王室仍可动员，却必须越来越仔细地核验谁能出役、何物能转运。南北压力并不自动合成一条通向覆亡的直线，但它们共同要求中心不断从王畿和封国取得人力、物资与协作。

\subsection{它留下了什么副作用}

王室以承认地方精英换取治理能力，也让地方精英把受命、资源、墓地和祖先记忆积累为跨代能力。王权强盛时，这种积累使网络更有弹性：封国能守通道、经营人群，也能在王室名分下彼此联系。王权难以持续赏赐、裁断和保护时，原先的协作渠道便可成为各方另行结盟的基础。这个反馈回路不是“分封必然导致分裂”的宿命论；它说明一种在有限行政条件下有效的制度，同时造就了中心日后必须面对的中介力量。

\eventref{WZ-0841-GONGHE}{西周·共和}已使王权与核心精英之间的裂缝公开化。到\eventref{WZ-0771-HAOJING}{西周·镐京陷落}，继承冲突、西部联盟和防御失调同时作用，关中的物质中心终于失守。东迁后，成周仍能承接宗庙、礼器和“天子”的语言，却无法把宗周的资源、军力和供给能力一同搬走。

\begin{turningpoint}[制度得失：衡量的不是一张制度标签]
\term{国家能力}：封国与双都使周能在交通、文书和常驻官员有限的条件下跨越较大空间；它的能力依赖中介者持续接命，不等于中心对地方的即时控制。

\term{财政与军事}：田宅、赏赐、仓储、铸铜和戍役让命令获得物质底座，也把动员成本分散到王畿和地方社会。器铭能证明资源与服务相连，不能替我们开出完整的税收或兵额账单。

\term{地方与文化}：共同的王命、礼器和祖先语言让不同地区可以互认；各地墓葬、工艺和资源的差异又始终存在。西周留下的不是一套不变的“周礼”，而是一种可被诸侯继续调用、也可被他们重新解释的政治语言。
\end{turningpoint}

这种遗产不会在前771年消失。东周诸侯仍要借周王室名分说话，春秋霸主仍在会盟中调用旧的礼仪词汇；只是越来越多能使词汇兑现的兵员、粮食、道路与地方网络，已掌握在诸侯手中。这正是下一章的起点。
